\chapter{CI/CD, containers, cloud systems, and infrastructure as code}

\section{Delivery is a controlled experiment}
Continuous integration (CI) makes changes build and test together. Continuous
delivery keeps a releasable artifact available; continuous deployment may
automatically release it after required checks. These are related practices, not
synonyms. The useful question is how quickly and safely the team can move from a
reviewed change to an observable, reversible production state.

A pipeline should answer:

\begin{enumerate}
  \item What exact source revision produced this artifact?
  \item Which tests, analyses, and security checks ran?
  \item Which configuration is injected at deployment rather than baked into the artifact?
  \item How is the artifact promoted, observed, rolled back, or rolled forward?
  \item Who or what is authorized to perform each transition?
\end{enumerate}

\section{Build once, configure at runtime}
Build a versioned artifact once and promote that artifact through environments.
Rebuilding separately for staging and production can produce different bytes and
invalidate test evidence. Configuration such as endpoints and feature flags may
vary; code and dependencies should not silently vary.

Pin dependencies and base images to a reviewed policy. “Latest” is convenient
for a tutorial and dangerous for reproducibility. Updates still need to happen;
the point is to make them explicit and observable.

\section{Containers}
A container packages a process and its user-space dependencies while sharing the
host kernel. It is not a complete security boundary by itself. Use a minimal
image, run as a non-root user, drop unnecessary capabilities, make the filesystem
read-only where feasible, bound resources, and avoid embedding secrets. The
example Compose file demonstrates several of these controls.

Docker's current documentation describes the Compose Specification as the
recommended format for defining services, networks, and volumes; the format is
implemented by modern Compose CLI versions \cite{docker2026compose}. Commands
and implementation details remain version-dependent, so CI should verify the
toolchain it expects.

\section{Infrastructure as code}
Infrastructure as code makes desired infrastructure reviewable and repeatable.
It does not make infrastructure deterministic automatically. Providers have
eventual consistency, external state, secrets, quotas, and destructive changes.

Treat infrastructure changes like application changes:

\begin{itemize}
  \item plan and review the diff;
  \item separate creation, migration, and deletion risk;
  \item protect state and limit its sensitive contents;
  \item test modules in a disposable environment;
  \item document manual recovery when automation is unavailable.
\end{itemize}

\section{Deployment strategies}
\noindent\begin{tabularx}{\textwidth}{@{}p{0.22\textwidth}p{0.27\textwidth}X@{}}
\toprule
Strategy & Benefit & Risk or requirement \\
\midrule
Rolling & Gradual replacement with one pool & Old and new versions must coexist \\
Blue-green & Fast switch between two environments & Doubles capacity and complicates state \\
Canary & Evidence from a small traffic slice & Requires good segmentation and rollback signals \\
Feature flag & Decouples code release from user exposure & Flag lifecycle and combinatorial states \\
\bottomrule
\end{tabularx}

Rollback is not always safe. A database migration, irreversible message, or
external side effect can make old code unable to run. Design forward-compatible
changes and a compensating action where possible.

\section{CI workflow example}
The repository workflow uses pinned major action versions, a fresh checkout,
separate example and PDF jobs, and an uploaded artifact. In a real project, add
permissions with least privilege, dependency review, test result retention,
protected environments, and secret handling appropriate to the platform.

\begin{exercise}
Design a release pipeline for a service with a schema migration and a queue
consumer. Mark which steps can be automatic, which require approval, and what
signal triggers rollback or a pause.
\end{exercise}

\section{Chapter summary}
Delivery should produce traceable, reproducible, observable transitions. Build
once, configure deliberately, harden containers, review infrastructure changes,
and choose deployment strategies based on compatibility and rollback reality.
